Let $a_{n,m} \coloneqq f(pw^{n+1}qx^m) - f(pw^{n}qx^m)$.
From Proposition~\ref{prop:convex}, we know $f(py^{n}qx^m)$ is $r$-convex in the variable $y$.
As a consequence, $f(pw^{n}qx^m)$ is convex in the variable $w$, i.e., 
\[
f(pw^{n+1}qx^m) + f(pw^{n-1}qx^m) \geq 2f(pw^{n}qx^m)
\]
which can be rewritten as
\[
f(pw^{n+1}qx^m) - f(pw^{n}qx^m) \geq f(pw^{n}qx^m) - f(pw^{n-1}qx^m)
\]
i.e., $a_{n,m}$ is increasing in $n$.
Let $b_{n,m} \coloneqq a_{n,m} - a_{n,m-1}$.
Adding and subtracting $f([x^m])$ and applying Proposition~\ref{cor:additivity-limit} twice, we get $\lim_{m} b_{n,m}$ exists.
Finally, we note that by Proposition~\ref{prop:bounded_double}, $b_{n,m} \leq 2f([wx])$.
Hence, the assumptions of Proposition~\ref{prop:exchange_limits} apply, and thus
\[
\lim_m \lim_n f(pw^{n+1}qx^m) - f(pw^{n}qx^m) = \lim_n \lim_m f(pw^{n+1}qx^m) - f(pw^{n}qx^m)
\]

Now, if we let $z=(p_1^{-1} \cdot \wt{x})^-$, and replace $y$ by $w=y^r$ in the above definition of $B_{a^{(n)}, b^{(n)}, c^{(n)}, d^{(n)}}$ we have
    \begin{multline*}
      \begin{aligned}
      \mu_f((\gamma_-,\gamma_+) \times [z,w \cdot z)) \\
    &= \lim_{n\to\infty} \mu_f(B_{a^{(n)}, b^{(n)}, c^{(n)}, d^{(n)}})  \\
    &= \frac{1}{2}\lim_{n\to\infty} \lim_{m\to\infty}f([pw^{n+1}qx^{m}]) + f([pw^{-(n+1)}qx^{m}])\\
    &-f([pw^{n}qx^{m}]) - f([pw^{-n}qx^{m}]) \\
    &= \frac{1}{2}\lim_{n\to\infty} \lim_{m\to\infty}f([pw^{n+1}qx^{m}]) -f([pw^{n}qx^{m}])  \\
     + \lim_{n\to\infty} \lim_{m\to\infty}  f([pw^{-(n+1)}qx^{m}]) - f([pw^{-n}qx^{m}])
    &= \frac{1}{2}\lim_{m\to\infty} \lim_{n\to\infty}f([pw^{n+1}qx^{m}]) -f([pw^{n}qx^{m}]) -  \\
      \lim_{m\to\infty} \lim_{n\to\infty}  f([pw^{-(n+1)}qx^{m}]) - f([pw^{-n}qx^{m}])
    &= \frac{1}{2}(f([w]) + f([w^{-1}]))\\
    &=f([w])
      \end{aligned}
    \end{multline*}      The first equality follows by the fact that the boxes defined above converge from below to $(\gamma_-,\gamma_+) \times [z,y
      \cdot z)$ as $n \to \infty$. The second equality is by
      definition of the measure $\mu_f$ on boxes.
      For the third, we split the nested limit into a sum of nested limits, for the fourth we exchange nested limits thanks to the above. For the fifth we add and subtract $f([y^{n+1}])$, and use the fact that $a_{n,m}$ is a monotone increasing and upper bounded sequence of $n$.
       The hypotheses of this Proposition are satisfied by the choice of $p_1$ and $q_1$.
      In the fourth equality, we then apply that
      Proposition to cancel out equivalent
      terms. In the sixth equality, we use the symmetry property of~$f$ to obtain the result.